\centerline{\bf \Huge Сириус}


\problem{1}
Найдите какое-нибудь натуральное число такое, что если из него вычесть удвоенную сумму его цифр, то получится число 2019.

\problem{2}
У Пети и Коли в стаканах налито равное количество молока. Вначале Петя перелил половину молока из своего стакана в стакан Коли, затем Коля перелил третью часть молока из своего стакана (с учетом молока, добавленного Петей) в Петин. Затем Петя перелил четверть молока из своего стакана в Колин, затем Коля перелил пятую часть и т.д. Сколько молока будет в стаканах у ребят после 2019-го переливания?

\problem{3}
На шахматной доске $8 \times 8$ стоят две фигуры: белая ладья на клетке $h1$ и черная ладья на клетке $h8$ . Сначала ход делает белая ладья, а затем черная. Сколько различных позиций может получиться? Позиции, отличающиеся поворотом доски, считаются различными, ладьи ходят по правилам шахмат.

\problem{4}
В выпуклом четырехугольнике ABCD угол A равен $60^{\circ}$. Точки M и N симметричны вершине А относительно прямых, содержащих стороны ВС и СD. Оказалось, что точки $\mathrm{B}, \mathrm{D}, \mathrm{M}, \mathrm{N}$ лежат на одной прямой. Найдите угол C четырехугольника ABCD.

\problem{5}
Можно ли подобрать четыре числа, не все из которых равны нулю так, чтобы их сумма квадратов была равна произведению одного из этих чисел на сумму остальных?

\problem{6}
У Васи и Пети есть по n монет. Каждая монета имеет достоинство 1 копейка, 50 копеек или 1 рубль. Оказалось, что у мальчиков поровну денег, но наборы монет не совпадают. Не обязательно, что монеты каждого достоинства присутствуют. При каком наименьшем n такое возможно?